% GENERATED FILE. Edit canonical lessons, then run npm run generate:books.
% canonical-source-hash: fnv1a64:93dc5b6682db2e94
% canonical-lessons: ZH-C486-zhaopin, ZH-C486-jianli, ZH-C486-yingpin, ZH-C486-shixi, ZH-C486-daiyu

\chapter{Looking for Work}
\label{ch:zh-looking-for-work}
\hlchaptermodality{486}

\begin{chapteropening}
\textbf{By the end of this chapter:} \emph{I can talk about job adverts, a CV, applying for a job, an internship and pay and conditions.}

Complete the last lesson of chapter 486: I can talk about job adverts, a CV, applying for a job, an internship and pay and conditions.
\end{chapteropening}

\section[zhāopìn]{zhāopìn — to recruit; to advertise a job}

\label{lesson:ZH-C486-zhaopin}



\begin{quote}
Before the new one: say the Mandarin for mindset, then the Mandarin for to relieve stress.
\end{quote}

\subsection*{You'll want to know: zhāopìn}
\emph{zhāopìn} — \textquotedblleft{}to recruit; to advertise a job\textquotedblright{}.

\begin{quote}
\emph{zhè jiā gōngsī zhèngzài zhāopìn} — This company is hiring.
\end{quote}

Its characters come later; for now it is a word for the ear.

\begin{cousinweb}
\emph{zhāopìn} is built from \emph{zhāo} 'to beckon' and \emph{pìn} 'to engage by invitation', both inherited from Middle Chinese, \emph{tsyew} and \emph{phjiengH} (Guangyun, in Baxter's transcription), which glosses the first 'to call to' and the second 'a formal visit of inquiry'; the sense 'to hire' is not given there.
\end{cousinweb}

\begin{grammarlens}[title={What you've built}]
One more word for saying how it was.
\end{grammarlens}

\subsection*{Your turn}
\begin{itemize}
\raggedright
  \item \emph{Say it:} \emph{zhāopìn}
  \item \emph{Say it:} \emph{zhāopìn}, once more
  \item \emph{Say it:} \emph{zhè jiā gōngsī zhèngzài zhāopìn} — This company is hiring.
  \item \emph{Recall:} say the Mandarin for mindset, then the Mandarin for to relieve stress, then say \emph{zhāopìn} again
\end{itemize}

\begin{tcolorbox}[breakable,colback=teal!4,colframe=teal!35!black,title={Before you move on}]
What does zhāopìn mean? (\textquotedblleft{}To recruit; to advertise a job\textquotedblright{}.) Say it once more.
\end{tcolorbox}
\section[jiǎnlì]{jiǎnlì — CV; résumé}

\label{lesson:ZH-C486-jianli}



\begin{quote}
Before the new one: say the Mandarin for to relieve stress, then the Mandarin for to recruit staff.
\end{quote}

\subsection*{You'll want to know: jiǎnlì}
\emph{jiǎnlì} — \textquotedblleft{}CV; résumé\textquotedblright{}.

\begin{quote}
\emph{zhè shì wǒ de jiǎnlì} — Here's my CV.
\end{quote}

Its characters come later; for now it is a word for the ear.

\begin{cousinweb}
\emph{jiǎnlì} is built from \emph{jiǎn} 'a writing slip; brief' and \emph{lì} 'to pass through', both inherited from Middle Chinese, \emph{keanX} and \emph{lek} (Guangyun, in Baxter's transcription), which glosses the first 'a writing tablet' and the second 'to pass through'. The sense 'a CV' is modern and not given there.
\end{cousinweb}

\begin{grammarlens}[title={What you've built}]
One more word for saying how it was.
\end{grammarlens}

\subsection*{Your turn}
\begin{itemize}
\raggedright
  \item \emph{Say it:} \emph{jiǎnlì}
  \item \emph{Say it:} \emph{jiǎnlì}, once more
  \item \emph{Say it:} \emph{zhè shì wǒ de jiǎnlì} — Here's my CV.
  \item \emph{Recall:} say the Mandarin for to relieve stress, then the Mandarin for to recruit staff, then say \emph{jiǎnlì} again
\end{itemize}

\begin{tcolorbox}[breakable,colback=teal!4,colframe=teal!35!black,title={Before you move on}]
What does jiǎnlì mean? (\textquotedblleft{}CV; résumé\textquotedblright{}.) Say it once more.
\end{tcolorbox}
\section[yìngpìn]{yìngpìn — to apply for a job}

\label{lesson:ZH-C486-yingpin}



\begin{quote}
Before the new one: say the Mandarin for to recruit staff, then the Mandarin for a CV.
\end{quote}

\subsection*{You'll want to know: yìngpìn}
\emph{yìngpìn} — \textquotedblleft{}to apply for a job\textquotedblright{}.

\begin{quote}
\emph{qù yìngpìn de shíhou xīntài zuì zhòngyào} — When you go to apply for a job, your frame of mind matters most.
\end{quote}

Its characters come later; for now it is a word for the ear.

\begin{cousinweb}
\emph{yìngpìn} is built from \zh{应} 'to answer' and \emph{pìn} 'to engage by invitation', both inherited from Middle Chinese, \emph{'ingH} and \emph{phjiengH} (Guangyun, in Baxter's transcription), which glosses the first 'things answering each other'.
\end{cousinweb}

\begin{grammarlens}[title={What you've built}]
One more word for saying how it was.
\end{grammarlens}

\subsection*{Your turn}
\begin{itemize}
\raggedright
  \item \emph{Say it:} \emph{yìngpìn}
  \item \emph{Say it:} \emph{yìngpìn}, once more
  \item \emph{Say it:} \emph{qù yìngpìn de shíhou xīntài zuì zhòngyào} — When you go to apply for a job, your frame of mind matters most.
  \item \emph{Recall:} say the Mandarin for to recruit staff, then the Mandarin for a CV, then say \emph{yìngpìn} again
\end{itemize}

\begin{tcolorbox}[breakable,colback=teal!4,colframe=teal!35!black,title={Before you move on}]
What does yìngpìn mean? (\textquotedblleft{}To apply for a job\textquotedblright{}.) Say it once more.
\end{tcolorbox}
\section[shíxí]{shíxí — to do an internship; to work as a trainee}

\label{lesson:ZH-C486-shixi}



\begin{quote}
Before the new one: say the Mandarin for a CV, then the Mandarin for to apply for a job.
\end{quote}

\subsection*{You'll want to know: shíxí}
\emph{shíxí} — \textquotedblleft{}to do an internship; to work as a trainee\textquotedblright{}.

\begin{quote}
\emph{wǒ zài yínháng shíxí le sān ge yuè} — I did a three-month internship at a bank.
\end{quote}

Its characters come later; for now it is a word for the ear.

\begin{cousinweb}
\emph{shíxí} is built from \emph{shí} 'real' and \emph{xí} 'to practise', both inherited from Middle Chinese, \emph{zyit} and \emph{zip} (Guangyun, in Baxter's transcription), which glosses the first 'full' and 'true' and the second 'to learn'.
\end{cousinweb}

\begin{grammarlens}[title={What you've built}]
One more word for saying how it was.
\end{grammarlens}

\subsection*{Your turn}
\begin{itemize}
\raggedright
  \item \emph{Say it:} \emph{shíxí}
  \item \emph{Say it:} \emph{shíxí}, once more
  \item \emph{Say it:} \emph{wǒ zài yínháng shíxí le sān ge yuè} — I did a three-month internship at a bank.
  \item \emph{Recall:} say the Mandarin for a CV, then the Mandarin for to apply for a job, then say \emph{shíxí} again
\end{itemize}

\begin{tcolorbox}[breakable,colback=teal!4,colframe=teal!35!black,title={Before you move on}]
What does shíxí mean? (\textquotedblleft{}To do an internship; to work as a trainee\textquotedblright{}.) Say it once more.
\end{tcolorbox}
\section[dàiyù]{dàiyù — pay and conditions (of a job)}

\label{lesson:ZH-C486-daiyu}



\begin{quote}
Before the new one: say the Mandarin for to apply for a job, then the Mandarin for to do an internship.
\end{quote}

\subsection*{You'll want to know: dàiyù}
\emph{dàiyù} — \textquotedblleft{}pay and conditions (of a job)\textquotedblright{}.

\begin{quote}
\emph{zhè jiā gōngsī de dàiyù bǐ nà jiā hǎo} — This company pays better than that one.
\end{quote}

Its characters come later; for now it is a word for the ear.

\begin{cousinweb}
\emph{dàiyù} is built from \emph{dài} 'to treat' and \emph{yù} 'to meet', both inherited from Middle Chinese, \emph{dojX} and \emph{ngjuH} (Guangyun, in Baxter's transcription), which glosses the first 'to wait for' and the second 'to meet without arranging it'; the sense 'how one is treated at work' is not given there.
\end{cousinweb}

\begin{grammarlens}[title={What you've built}]
One more word for saying how it was.
\end{grammarlens}

\subsection*{Your turn}
\begin{itemize}
\raggedright
  \item \emph{Say it:} \emph{dàiyù}
  \item \emph{Say it:} \emph{dàiyù}, once more
  \item \emph{Say it:} \emph{zhè jiā gōngsī de dàiyù bǐ nà jiā hǎo} — This company pays better than that one.
  \item \emph{Recall:} say the Mandarin for to apply for a job, then the Mandarin for to do an internship, then say \emph{dàiyù} again
\end{itemize}

\begin{tcolorbox}[breakable,colback=teal!4,colframe=teal!35!black,title={Before you move on}]
What does dàiyù mean? (\textquotedblleft{}Pay and conditions (of a job)\textquotedblright{}.) Say it once more.
\end{tcolorbox}
